% \documentclass[aspectratio=169,notes]{beamer}
\documentclass[aspectratio=169]{beamer}
\usetheme[faculty=phil]{fibeamer}
\usepackage{polyglossia}
\setmainlanguage{english} %% main locale instead of `english`, you
%% can typeset the presentation in either Czech or Slovak,
%% respectively.
\setotherlanguages{russian} %% The additional keys allow
%%
%%   \begin{otherlanguage}{czech}   ... \end{otherlanguage}
%%   \begin{otherlanguage}{slovak}  ... \end{otherlanguage}
%%
%% These macros specify information about the presentation
\title[Artificial Intelligence Software (AIS)]{Artificial Intelligence Software, Lecture 7} %% that will be typeset on the
\subtitle{Training process
\\ \ \\ \  } %% title page.
\author{Oleg Bulichev}
%% These additional packages are used within the document:
\usepackage{ragged2e}  % `\justifying` text
\usepackage{booktabs}  % Tables
\usepackage{tabularx}
\usepackage{tikz}      % Diagrams
\usetikzlibrary{calc, shapes, backgrounds}
\usetikzlibrary{decorations.pathreplacing,calligraphy,calc,graphs}
\usepackage{amsmath, amssymb}
\usepackage{url}       % `\url`s
\usepackage{listings}  % Code listings
\usepackage{subfigure}
% \usepackage{floatrow}
% \usepackage{subcaption}
\usepackage{mathtools}
\usepackage{todonotes}
\usepackage{fontspec}
\usepackage{multicol}
\usepackage{pdfpages}
\usepackage{wrapfig}
\usepackage{qrcode}
\usepackage{animate}
\usepackage{booktabs}
\usepackage{multirow}
\usepackage{graphicx}
\usepackage{colortbl}

\graphicspath{{resources/}}
\frenchspacing

\setbeamertemplate{caption}{\raggedright\insertcaption\par}
\usetikzlibrary{graphs}
\usetikzlibrary{positioning,arrows.meta}

% \usepackage[backend=biber,style=ieee,autocite=footnote]{biblatex}
% \addbibresource{biblio.bib}
% \DefineBibliographyStrings{english}{%
%   bibliography = {References},}

\newcommand{\oleg}[2][] {\todo[color=red, #1] {OLEG:\\ #2}}
\newcommand{\fbckg}[1]{\usebackgroundtemplate{\includegraphics[width=\paperwidth]{#1}}}%frame background

\usepackage[framemethod=TikZ]{mdframed}
\newcommand{\dbox}[1]{
\begin{mdframed}[roundcorner=3pt, backgroundcolor=yellow, linewidth=0]
\vspace{1mm}
{#1}
\vspace{1mm}
\end{mdframed}
}

\newcommand{\imgmock}[2]{
\fbox{\parbox[c][0.48\textheight][c]{0.92\linewidth}{\centering\textbf{Mock image placeholder}\\[2mm]#1\\[2mm]\footnotesize #2}}
}

\begin{document}
\setlength{\abovedisplayskip}{0pt}
\setlength{\belowdisplayskip}{0pt}
\setlength{\abovedisplayshortskip}{0pt}
\setlength{\belowdisplayshortskip}{0pt}

\fbckg{fibeamer/figs/title_page.png}
\frame[c]{\setcounter{framenumber}{0}
    \usebeamerfont{title}%
    \usebeamercolor[fg]{title}%
    \begin{minipage}[b][6.5\baselineskip][b]{\textwidth}%
        \textcolor{black}{\raggedright\inserttitle}
    \end{minipage}
    % \vskip-1.5\baselineskip

    \usebeamerfont{subtitle}%
    \usebeamercolor[fg]{framesubtitle}%
    \begin{minipage}[b][3\baselineskip][b]{\textwidth}
        \raggedright%
        \insertsubtitle%
    \end{minipage}
    \vskip.25\baselineskip
}
%   \frame[c]{\maketitle}

\fbckg{fibeamer/figs/common.png}

\note{\scriptsize \begin{itemize}
        \item \
    \end{itemize}}


\begin{frame}[t]{Training as a Pipeline}
Training a neural network is a production pipeline, not one monolithic algorithm.
\vspace{1mm}
\begin{itemize}
    \item Input stage: data loading, filtering, augmentation, batching.
    \item Compute stage: forward pass, loss computation, backward pass, optimizer step.
    \item Control stage: validation, checkpointing, logging, early stopping.
    \item Ops stage: experiment tracking, distributed execution, model registry.
\end{itemize}
\vspace{1mm}
Your work is to design and control this system, not to re-implement autograd.
\end{frame}

\begin{frame}{Pipeline Scheme}
\centering
\begin{tikzpicture}[
    stage/.style={draw, rounded corners, align=center, minimum width=2.0cm, minimum height=0.95cm, fill=blue!8},
    arrow/.style={-{Latex[length=2.5mm]}, thick}
]
\node[stage] (d) {Dataset\\+ transforms};
\node[stage, right=0.8cm of d] (l) {DataLoader\\(batches)};
\node[stage, right=0.8cm of l] (f) {Forward\\pass};
\node[stage, right=0.8cm of f] (loss) {Loss};
\node[stage, right=0.8cm of loss] (b) {Backward\\+ optimizer};

\draw[arrow] (d) -- (l);
\draw[arrow] (l) -- (f);
\draw[arrow] (f) -- (loss);
\draw[arrow] (loss) -- (b);
\draw[arrow, bend left=-20] (b.north) to node[above, font=\scriptsize] {next batch} (l.north);

\node[draw, rounded corners, fill=green!10, below=0.95cm of f, minimum width=3.0cm, minimum height=0.8cm] (val) {Valid. loop};
\node[draw, rounded corners, fill=yellow!15, below=0.95cm of loss, minimum width=3.0cm, minimum height=0.8cm] (track) {Tracking + checkpoints};
\draw[arrow] (f.south) -- (val.north);
\draw[arrow] (loss.south) -- (track.north);
\end{tikzpicture}
\end{frame}

\begin{frame}[t]{Key Idea: Build on Existing Models}
In practical ML, we rarely train foundation models from scratch.
\vspace{1mm}
\begin{itemize}
    \item Start with pretrained encoders/LLMs/vision backbones.
    \item Adapt with fine-tuning, adapters, or LoRA-style parameter-efficient methods.
    \item Keep a consistent pipeline: same data contracts, metrics, and deployment checks.
\end{itemize}
\vspace{1mm}
Engineering maturity comes from reproducibility and iteration speed.
\end{frame}

\begin{frame}[t]{Fine-Tuning Strategies}
\begin{itemize}
    \item \textbf{Full fine-tuning:} best quality ceiling, highest memory/compute cost.
    \item \textbf{Partial unfreezing:} train only top layers for faster adaptation.
    \item \textbf{PEFT (LoRA/adapters):} small trainable modules, cheap and scalable.
\end{itemize}
\vspace{1mm}
Typical workflow:
\begin{enumerate}
    \item Freeze most layers.
    \item Train head or adapters.
    \item Unfreeze gradually if validation metric plateaus.
\end{enumerate}
\end{frame}

\begin{frame}[t]{Data and Dataset Quality}
Dataset quality determines your quality ceiling.
\vspace{1mm}
\begin{itemize}
    \item Remove duplicates and label leakage between train/validation/test.
    \item Control class imbalance and long-tail behavior.
    \item Keep metadata for slicing metrics (domain, source, language, device).
    \item Version data exactly as you version code.
\end{itemize}
\vspace{1mm}
Bad labels are often a bigger bottleneck than model architecture.
\end{frame}

\begin{frame}[t]{DataLoader Design Decisions}
DataLoader converts storage format into GPU-ready batches.
\vspace{1mm}
\begin{itemize}
    \item Throughput: prefetch, pinned memory, parallel workers.
    \item Stochasticity: shuffling and sampling strategy.
    \item Consistency: deterministic seeds for reproducible debugging.
    \item Efficiency: dynamic padding or bucketing for variable-length inputs.
\end{itemize}
\end{frame}

\begin{frame}[t]{Batching and Effective Batch Size}
\begin{itemize}
    \item Physical batch size is limited by GPU memory.
    \item Effective batch size can be increased via gradient accumulation.
    \item Very large batches may require learning-rate scaling and warmup.
\end{itemize}
\vspace{1mm}
\[
\text{effective batch} = \text{batch per device} \times \text{num devices} \times \text{accumulation steps}
\]
\end{frame}

\begin{frame}[t]{Forward Pass and Loss}
\begin{itemize}
    \item Forward pass transforms batch input into logits or token probabilities.
    \item Loss encodes the optimization objective (cross-entropy, MSE, contrastive, etc.).
    \item Metrics are not always losses: optimize one thing, report several.
\end{itemize}
\vspace{1mm}
A common anti-pattern is optimizing a proxy loss that does not match business KPI.
\end{frame}

\begin{frame}[t]{Backward Pass and Optimization}
\begin{itemize}
    \item Autograd computes gradients automatically from loss to parameters.
    \item Optimizers (AdamW, SGD) update weights using gradients.
    \item Schedulers adjust learning rate over time (warmup, cosine decay, step decay).
    \item Mixed precision speeds training and reduces memory footprint.
\end{itemize}
\end{frame}

\begin{frame}[t]{Validation and Error Analysis}
Validation should answer: \textit{is the model good enough for deployment?}
\vspace{1mm}
\begin{itemize}
    \item Track both aggregate metrics and slices (rare classes, edge domains).
    \item Compare calibration, confidence, and failure modes.
    \item Keep a fixed benchmark subset for stable trend monitoring.
\end{itemize}
\vspace{1mm}
Training loss without structured error analysis is misleading.
\end{frame}

\begin{frame}[t]{Experiment Tracking and Reproducibility}
\begin{itemize}
    \item Log config, git revision, dataset version, metrics, artifacts.
    \item Save checkpoints and best model according to validation target.
    \item Use systems such as ClearML/MLflow/W\&B to compare runs.
\end{itemize}
\vspace{1mm}
Reproducibility means you can recreate the same run months later.
\end{frame}

\begin{frame}[t]{Hyperparameters That Matter Most}
\begin{tabular}{@{}p{3.2cm}p{8.0cm}@{}}
\toprule
Group & Examples and effect \\
\midrule
Optimization & Learning rate, weight decay, optimizer type; strongest impact on convergence \\
Batching & Batch size, accumulation steps; trade-off between stability and speed \\
Regularization & Dropout, label smoothing, augmentation; controls overfitting \\
Schedule & Warmup ratio, max epochs, early stopping; affects final generalization \\
\bottomrule
\end{tabular}
\end{frame}

\begin{frame}[t]{Why Ray for Training Workloads}
Ray is a distributed execution framework for Python workloads.
\vspace{1mm}
\begin{itemize}
    \item Scale one script from laptop to cluster with minimal code changes.
    \item Ray Train: distributed model training (data parallel and fault-tolerant loops).
    \item Ray Tune: large-scale hyperparameter search and scheduler-based pruning.
    \item Ray Data: scalable data ingestion and preprocessing pipeline.
\end{itemize}
\vspace{1mm}
Ray is useful when experiments outgrow a single machine.
\end{frame}

\begin{frame}[t]{Ray in the Training Stack}
\centering
\begin{tikzpicture}[
    box/.style={draw, rounded corners, align=center, minimum width=3.4cm, minimum height=0.95cm, fill=orange!12},
    arr/.style={-{Latex[length=2.5mm]}, thick}
]
\node[box] (user) {Training script\\(PyTorch/Lightning)};
\node[box, below=0.7cm of user] (train) {Ray Train\\distributed workers};
\node[box, left=1.2cm of train] (data) {Ray Data\\input pipeline};
\node[box, right=1.2cm of train] (tune) {Ray Tune\\HPO + schedulers};
\node[box, below=0.7cm of train, minimum width=5.2cm] (cluster) {Ray Cluster\\head node + worker nodes + autoscaling};

\draw[arr] (user) -- (train);
\draw[arr] (data) -- (train);
\draw[arr] (tune) -- (train);
\draw[arr] (train) -- (cluster);
\end{tikzpicture}
\end{frame}

\begin{frame}[t]{Ray Concepts to Remember}
\begin{itemize}
    \item \textbf{Tasks:} parallel stateless functions.
    \item \textbf{Actors:} stateful workers for long-lived services.
    \item \textbf{Object Store:} shared-memory transport for large tensors and batches.
    \item \textbf{Schedulers:} resource-aware placement (CPU/GPU/custom resources).
\end{itemize}
\vspace{1mm}
These abstractions let you scale preprocessing, training, and HPO consistently.
\end{frame}

\begin{frame}[t]{Ray Tune for Hyperparameter Search}
\begin{itemize}
    \item Define search space (learning rate, batch size, dropout, optimizer).
    \item Run many trials in parallel across cluster resources.
    \item Use pruning schedulers (e.g., ASHA) to stop weak trials early.
    \item Keep top checkpoints and export best configuration.
\end{itemize}
\vspace{1mm}
Result: faster iteration and better utilization of available GPUs.
\end{frame}

\begin{frame}[t]{PyTorch Distributed: What and When}
\begin{itemize}
    \item \textbf{torch.distributed (DDP)} is the default baseline for multi-node training in both industry and academia.
    \item Most common CPU backend: \texttt{gloo}; for GPUs in production: \texttt{nccl}.
    \item Typical launch pattern: \texttt{torchrun} + one process per device (or one per CPU node in demos).
    \item Best fit: stable training jobs where you control infrastructure and need maximum performance.
\end{itemize}
\vspace{1mm}
Rule of thumb: start from DDP mental model first, then introduce Ray for orchestration and scale management.
\end{frame}

\begin{frame}[t]{DDP vs Ray Train: Practical Choice}
\begin{tabular}{@{}p{3.0cm}p{4.0cm}p{4.3cm}@{}}
\toprule
Aspect & PyTorch Distributed & Ray Train \\
\midrule
Abstraction & Low-level, explicit control & Higher-level orchestration \\
Setup & Manual launch and env setup & Unified cluster/runtime layer \\
Fault handling & Mostly manual & Built-in recovery patterns \\
Best for & Peak performance, custom loops & Many experiments, elastic scaling \\
\bottomrule
\end{tabular}
\vspace{1mm}
In practice teams combine both: DDP-style training logic wrapped by Ray execution.
\end{frame}

\begin{frame}[t,fragile]{CPU Cluster Emulation with Docker Compose}
\centering
\begin{tikzpicture}[
    nodebox/.style={draw, rounded corners, align=center, minimum width=3.3cm, minimum height=0.95cm, fill=cyan!12},
    arr/.style={-{Latex[length=2.5mm]}, thick}
]
\node[nodebox, fill=orange!15] (compose) {docker-compose\\control plane};
\node[nodebox, below left=1.2cm and 1.2cm of compose] (master) {torch-node-1\\master rank 0};
\node[nodebox, below=1.2cm of compose] (worker1) {torch-node-2\\worker rank 1};
\node[nodebox, below right=1.2cm and 1.2cm of compose] (worker2) {torch-node-N\\worker rank N-1};
\node[nodebox, below=1.1cm of worker1, fill=green!14, minimum width=6.0cm] (ray) {ray-head + ray-worker (scaled)};

\draw[arr] (compose) -- (master);
\draw[arr] (compose) -- (worker1);
\draw[arr] (compose) -- (worker2);
\draw[arr] (compose) -- (ray);
\end{tikzpicture}
\vspace{1mm}
\begin{itemize}
    \item \texttt{WORLD\_SIZE=3 docker compose ... --scale torch-node=3}
    \item \texttt{docker compose ... --scale ray-worker=3}
\end{itemize}
\end{frame}

\begin{frame}[t]{When to Show Demo in This Lecture}
Recommended pacing for a 90-minute class:
\begin{enumerate}
    \item \textbf{After pipeline basics (minute 20--25):} run quick YOLO CPU fine-tune (1 epoch).
    \item \textbf{After reproducibility section (minute 40--45):} show ClearML run page and tracked artifacts.
    \item \textbf{After distributed theory (minute 60--70):} run Docker-based DDP and then Ray worker scaling.
\end{enumerate}
\vspace{1mm}
Keep each live demo under 3--5 minutes and prepare screenshots as fallback.
\end{frame}

\begin{frame}[t]{Mock Slide: What to Insert (System Architecture)}
\imgmock{Target image: \textbf{"End-to-end ML training architecture"}}{Show data source, preprocessing, feature store, training workers, validation service, model registry, and deployment endpoint. Add arrows for offline training flow and online inference flow in different colors.}
\end{frame}

\begin{frame}[t]{Mock Slide: What to Insert (Ray Cluster)}
\imgmock{Target image: \textbf{"Ray cluster topology for training"}}{Draw one head node and 3--6 worker nodes with CPU/GPU labels. Visualize task queue, object store, and autoscaler. Add a small legend: blue for data, orange for training, green for tuning trials.}
\end{frame}

\begin{frame}[t]{Mock Slide: What to Insert (Live Demo Setup)}
\imgmock{Target image: \textbf{"Classroom live-demo command map"}}{Create a terminal-style infographic with 3 blocks: (1) YOLO fine-tune command, (2) ClearML dashboard screenshot area, (3) Docker scale commands for torch-node and ray-worker. Add expected runtime labels: 2--5 min each.}
\end{frame}

\begin{frame}[t]{Tools Ecosystem}
\begin{itemize}
    \item \textbf{Modeling:} PyTorch, Lightning, Hugging Face Transformers.
    \item \textbf{Tracking:} ClearML (primary) or MLflow for metrics, configs, and artifacts.
    \item \textbf{Scaling:} Ray for distributed data processing, training, and HPO.
    \item \textbf{Serving:} model registry + deployment pipeline with rollback strategy.
\end{itemize}
\end{frame}

\begin{frame}[t]{Conclusion}
Training is engineering orchestration, not only gradient updates.
\vspace{1mm}
\begin{itemize}
    \item Understand each pipeline component and its failure modes.
    \item Optimize data quality and reproducibility before model complexity.
    \item Use Ray when scale and experiment throughput become bottlenecks.
\end{itemize}
\end{frame}


\fbckg{fibeamer/figs/last_page.png}
\frame[plain]{}

\end{document}
